\begin{center}
    {\LARGE \textbf{2022无机化学与物理化学}} \\[2ex]
    {\large 时量180分钟 \quad 满分150分}
\end{center}

\vspace{0.5cm}
\noindent\textbf{注意事项:}
\begin{enumerate}[label=\arabic*., parsep=0pt, itemsep=0pt]
    \item 答题前,考生先将自己的姓名、学号、考场号、座位号填写在答题卡上,将条形码准确粘贴在考生信息条形码粘贴区。
    \item 考试结束后,将本试卷与答题纸一并收回。
    \item 回答各个问题时,将答案写在答题纸上,写在本试卷上无效。
\end{enumerate}

\vspace{1cm}
\begin{center}
    {\Large \textbf{无机化学部分(共75分)}}
\end{center}

\noindent\textbf{一、选择题(20 pts.)}

\begin{enumerate}[label=\arabic*., itemsep=1.5ex]
    \item 下列配合物中，最不稳定的是
    \begin{tasks}(4)
        \task $[\ce{Cu(NH3)4}]^{2+}$
        \task $[\ce{Cu(H2O)4}]^{2+}$
        \task $[\ce{Cu(CN)4}]^{2-}$
        \task $[\ce{Cu(en)2}]^{2+}$
    \end{tasks}

    \item 下列物质中，不能将 $\ce{HI}$ 氧化的是
    \begin{tasks}(4)
        \task $\ce{Fe2O3}$
        \task $\ce{ZnO}$
        \task $\ce{MnO2}$
        \task $\ce{CuO}$
    \end{tasks}

    \item 下列化合物中，还原能力最强的是
    \begin{tasks}(4)
        \task $\ce{H2SO3}$
        \task $\ce{H2O}$
        \task $\ce{HI}$
        \task $\ce{HNO2}$
    \end{tasks}

    \item 相同浓度下的下列溶液中，碱性最强的是
    \begin{tasks}(4)
        \task $\ce{NaHSO3}$
        \task $\ce{NaHS}$
        \task $\ce{NaHSO4}$
        \task $\ce{NaHCO3}$
    \end{tasks}

    \item 下列各对元素中，性质最相似的是
    \begin{tasks}(4)
        \task $\ce{V}$ 和 $\ce{Ta}$
        \task $\ce{Ti}$ 和 $\ce{Zr}$
        \task $\ce{Eu}$ 和 $\ce{Gd}$
        \task $\ce{Fe}$ 和 $\ce{Ni}$
    \end{tasks}

    \item 下列分子中，具有大 $\pi$ 键的是
    \begin{tasks}(4)
        \task $\ce{SO2Cl2}$
        \task $\ce{CO2}$
        \task $\ce{SOCl2}$
        \task $\ce{HPO3}$
    \end{tasks}

    \item 下列配合物中，没有旋光异构体的是
    \begin{tasks}(2)
        \task $\ce{K3[Fe(C2O4)3]}$
        \task $[\ce{Ni(en)3}]Cl2$
        \task $[\ce{Pt(NH3)2Cl3(OH)}]$
        \task $[\ce{Co(en)2(NH3)2}]Cl2$
    \end{tasks}

    \item 下列硫化物中，能溶于 $\ce{Na2S}$ 溶液的是
    \begin{tasks}(4)
        \task $\ce{SnS}$
        \task $\ce{SnS2}$
        \task $\ce{PbS}$
        \task $\ce{Bi2S3}$
    \end{tasks}

    \item 下列氢化物中，最稳定的是
    \begin{tasks}(4)
        \task $\ce{PH3}$
        \task $\ce{AsH3}$
        \task $\ce{SbH3}$
        \task $\ce{BiH3}$
    \end{tasks}

    \item 下列氢氧化物中，不能被空气氧化的是
    \begin{tasks}(4)
        \task $\ce{Fe(OH)2}$
        \task $\ce{Co(OH)2}$
        \task $\ce{Ni(OH)2}$
        \task $\ce{Mn(OH)2}$
    \end{tasks}
\end{enumerate}

\vspace{1em}
\noindent\textbf{二、简要回答下列问题(45 pts.)}

\begin{enumerate}[label=\arabic*., itemsep=1.5ex]
    \item 解释实验现象并写出相关的化学反应方程式 (15 pts.)
    \begin{enumerate}[label=(\arabic*), noitemsep, topsep=0pt, partopsep=0pt]
        \item 向 $\ce{CuCl}$ 中滴加氨水，白色沉淀溶解成无色溶液，后变深蓝色。
        \item 少量硫代硫酸钠溶液和硝酸银溶液反应，先生成白色沉淀，沉淀逐渐变黄变棕最后变黑。
        \item $\ce{KI}$ 与 $\ce{NaNO2}$ 混合溶液加入稀硫酸后溶液立即变黄，$\ce{KI}$ 与 $\ce{NaNO3}$ 混合溶液加入稀硫酸后溶液不变色。
    \end{enumerate}

    \item 试解释：F 的电子亲合能比 Cl 的电子亲合能小，但 $\ce{F2}$ 反而比 $\ce{Cl2}$ 活泼。

    \item 已知 $[\ce{Co(NO2)6}]^{4-}$ 的磁矩为 1.8 B.M.，试讨论中心离子的杂化方式并预测其稳定性。

    \item 将 $\ce{NaNO2}$ 酸化后滴加 $\ce{KI}$ 溶液或将 $\ce{KI}$ 酸化后滴加 $\ce{NaNO2}$ 溶液都可制备 $\ce{NO}$。哪种制备方法更好？为什么？

    \item 设计方案将溶液中的 $\ce{Ag+}$、$\ce{Al^{3+}}$、$\ce{Cr^{3+}}$、$\ce{Mn^{2+}}$ 分离。

    \item 以碳酸钠和硫黄为原料制备硫代硫酸钠。

    \item 试用 3 种方法鉴别 $\ce{MnO2}$ 和 $\ce{PbO2}$，给出实验现象和相应的化学反应方程式。
\end{enumerate}

\vspace{1em}
\noindent\textbf{三、推断题(10 pts.)}

红色固体 A 溶于热的浓盐酸放出气体 B，溶液经冷却后析出绿色晶体 C。少量 B 通入 $\ce{KI}$ 溶液中则溶液变黄。向 C 的水溶液中加入碳酸钠溶液生成沉淀 D。将 D 溶于 $\ce{KOH}$ 溶液后加入 $\ce{H2O2}$ 溶液得到黄色溶液 E。用稀硫酸酸化溶液 E 后，经蒸发、浓缩、冷却有橘红色晶体 F 析出。F 与浓硫酸作用能够得到 A。给出 A、C、D、E 所代表的物质的化学式，并写出 A→B、D→E 和 F→A 各过程的方程式。

\vspace{0.5cm}
\begin{center}
    {\Large \textbf{物理化学部分(共75分)}}
\end{center}

\vspace{0.5cm}
\noindent\textbf{四、选择题(24 pts. total, 2 pts. each)}

\begin{enumerate}[label=\arabic*., start=1, itemsep=1.5ex]
    \item 下述说法哪一种不正确？
    \begin{tasks}(1) 
        \task 理想气体经绝热自由膨胀后，其内能变化为零
        \task 非理想气体经绝热自由膨胀后，其内能变化不一定为零
        \task 非理想气体经绝热膨胀后，其温度一定降低
        \task 非理想气体经一不可逆循环，其内能变化为零
    \end{tasks}

    \item 下述方法中，哪一种对于消灭蚂蟥比较有效？
    \begin{tasks}(4)
        \task 击打
        \task 刀割
        \task 晾晒
        \task 撒盐
    \end{tasks}

    \item 在碰撞理论中校正因子 $P$ 小于 1 的主要因素是：
    \begin{tasks}(2)
        \task 反应体系是非理想的
        \task 空间的位阻效应
        \task 分子碰撞的激烈程度不够
        \task 分子间的作用力
    \end{tasks}

    \item 298 K 时两个级数相同的反应 I、II，活化能 $E_\text{I} = E_\text{II}$，若速率常数 $k_\text{I} = 10 k_\text{II}$，则两反应之活化熵相差：
    \begin{tasks}(2)
        \task $0.6\ \text{J}\cdot\text{K}^{-1}\cdot\text{mol}^{-1}$
        \task $10\ \text{J}\cdot\text{K}^{-1}\cdot\text{mol}^{-1}$
        \task $19\ \text{J}\cdot\text{K}^{-1}\cdot\text{mol}^{-1}$
        \task $190\ \text{J}\cdot\text{K}^{-1}\cdot\text{mol}^{-1}$
    \end{tasks}

    \item 一恒压反应体系，若产物与反应物的 $\Delta C_p > 0$，则此反应
    \begin{tasks}(2)
        \task 吸热
        \task 放热
        \task 无热效应
        \task 吸放热不能肯定
    \end{tasks}

    \item 理想气体反应 $\ce{CO(g) + 2H2(g) = CH3OH(g)}$ 的 $\Delta_\text{r} G_\text{m}^\ominus$ 与温度 $T$ 的关系为：\\
    $\Delta_\text{r} G_\text{m}^\ominus / \text{J}\cdot\text{mol}^{-1} = -21660 + 52.92(T/\text{K})$，若使在标准状态下的反应向右进行，则应控制反应的温度：
    \begin{tasks}(2)
        \task 必须高于 409.3 K
        \task 必须低于 409.3 K
        \task 必须等于 409.3 K
        \task 必须低于 $409.3\ ^\circ\text{C}$
    \end{tasks}

    \item 金属活泼性排在 $\ce{H2}$ 之前的金属离子，如 $\ce{Na+}$ 能优先于 $\ce{H+}$ 在汞阴极上析出，这是由于：（ ）
    \begin{tasks}(1)
        \task $\phi^\ominus(\ce{Na+ / Na}) < \phi^\ominus(\ce{H+ / H2})$
        \task $\eta(\ce{Na}) < \eta(\ce{H2})$
        \task $\phi(\ce{Na+ / Na}) < \phi(\ce{H+ / H2})$
        \task $\ce{H2}$ 在汞上析出有很大的超电势，以至于 $\phi(\ce{Na+ / Na}) > \phi(\ce{H+ / H2})$
    \end{tasks}

    \item 使用瑞利（Reyleigh）散射光强度公式，在下列问题中可以解决的问题是
    \begin{tasks}(2)
        \task 溶胶粒子的大小
        \task 溶胶粒子的形状
        \task 测量散射光的波长
        \task 测量散射光的振幅
    \end{tasks}

    \item 一定量的理想气体从同一始态出发，分别经 (1)等温压缩，(2)绝热压缩到具有相同压力的终态，以 $H_1$、$H_2$ 分别表示两个终态的焓值，则有：
    \begin{tasks}(4)
        \task $H_1 > H_2$
        \task $H_1 = H_2$
        \task $H_1 < H_2$
        \task $H_1 \le H_2$
    \end{tasks}

    \item 在还原性酸性溶液中，$\ce{Zn}$ 的腐蚀速度较 $\ce{Fe}$ 为小，其原因是
    \begin{tasks}(1)
        \task $\phi(\ce{Zn^{2+} / Zn})(\text{平}) < \phi(\ce{Fe^{2+} / Fe})(\text{平})$
        \task $\phi(\ce{Zn^{2+} / Zn}) < \phi(\ce{Fe^{2+} / Fe})$
        \task $\phi(\ce{H+ / H2})(\text{平, } \ce{Zn}) < \phi(\ce{H+ / H2})(\text{平, } \ce{Fe})$
        \task $\phi(\ce{H+ / H2})(\ce{Zn}) < \phi(\ce{H+ / H2})(\ce{Fe})$
    \end{tasks}

    \item 已知有下列一组公式可用于理想气体绝热过程功的计算：
    

    \vspace{1ex}
    \begin{tabular}{p{0.45\textwidth} p{0.45\textwidth}}
        (1) $W = C_V(T_1 - T_2)$ & (2) $[1/(\gamma - 1)](p_1 V_1 - p_2 V_2)$ \\[1ex]
        (3) $[p_1 V_1 / (\gamma - 1)][1 - (V_1/V_2)^{\gamma - 1}]$ & (4) $[p_1 V_1 / (\gamma - 1)][1 - (p_2/p_1)^{1 - 1/\gamma}]$ \\[1ex]
        (5) $[p_1 V_1 / (\gamma - 1)][1 - (p_2 V_2 / p_1 V_1)]$ & (6) $[R/(\gamma - 1)](T_1 - T_2)$
    \end{tabular}
    \vspace{1ex}
    
    但这些公式只适于绝热可逆过程的是：
    \begin{tasks}(4)
        \task (1), (2)
        \task (3), (4)
        \task (5), (6)
        \task (4), (5)
    \end{tasks}

    \item 以石墨为阳极，电解 $0.01\ \text{mol}\cdot\text{kg}^{-1}\ \ce{NaCl}$ 溶液，在阳极上首先析出：
    \begin{tasks}(2)
        \task $\ce{Cl2}$
        \task $\ce{O2}$
        \task $\ce{Cl2}$ 与 $\ce{O2}$ 混合气体
        \task 无气体析出
    \end{tasks}
    已知：$\phi^\ominus(\ce{Cl2/Cl-}) = 1.36\ \text{V}$，$\eta(\ce{Cl2}) = 0\ \text{V}$，\\
    \hspace*{3em}$\phi^\ominus(\ce{O2/OH-}) = 0.401\ \text{V}$，$\eta(\ce{O2}) = 0.8\ \text{V}$。
\end{enumerate}

\vspace{0.5cm}
\noindent\textbf{五、计算题(10 pts.)}

$1.22 \times 10^{-2}\ \text{kg}$ 苯甲酸溶于 $0.10\ \text{kg}$ 乙醇后，使乙醇的沸点升高了 $1.13\ \text{K}$。若将 $1.22 \times 10^{-2}\ \text{kg}$ 苯甲酸溶于 $0.10\ \text{kg}$ 苯中，则苯的沸点升高 $1.36\ \text{K}$。计算苯甲酸在两种溶剂中的摩尔质量。计算结果说明什么问题。已知：$K_b(\text{乙醇}) = 1.19\ \text{K}\cdot\text{kg}\cdot\text{mol}^{-1}$，$K_b(\text{苯}) = 2.60\ \text{K}\cdot\text{kg}\cdot\text{mol}^{-1}$。

\vspace{0.5cm}
\noindent\textbf{六、计算题(15 pts.)}

\begin{enumerate}[label=（\arabic*）, itemsep=1ex]
    \item 利用下面数据绘制 Mg-Cu 体系的相图。Mg（熔点为 $648\ ^\circ\text{C}$）和 Cu（熔点为 $1085\ ^\circ\text{C}$）形成两种化合物：$\ce{MgCu2}$（熔点为 $800\ ^\circ\text{C}$），$\ce{Mg2Cu}$（熔点为 $580\ ^\circ\text{C}$）；形成三个低共熔混合物，其组成分别为 $w(\text{Mg})=0.10, w(\text{Mg})=0.33, w(\text{Mg})=0.65$，它们的熔点分别为 $690\ ^\circ\text{C}, 560\ ^\circ\text{C}$ 和 $380\ ^\circ\text{C}$。已知 $M_r(\text{Cu})=63.55\ \text{g}\cdot\text{mol}^{-1}, M_r(\text{Mg})=24.31\ \text{g}\cdot\text{mol}^{-1}$
    \item 画出组成为 $w(\text{Mg})=0.25$ 的 $900\ ^\circ\text{C}$ 的熔体降温到 $100\ ^\circ\text{C}$ 时的步冷曲线。
    \item 当(2)中的熔体 $1\ \text{kg}$ 降温至无限接近 $560\ ^\circ\text{C}$ 时，可得到何种化合物？其质量为若干？
\end{enumerate}

\vspace{0.5cm}
\noindent\textbf{七、计算题(13 pts.)}

电池 $\ce{Ag(s) | AgBr(s) | HBr(0.1\ mol\cdot kg^{-1}) | H2(0.01\ p^\ominus), Pt}$，$298\ \text{K}$ 时，$E = 0.165\ \text{V}$。当电子得失为 $1\ \text{mol}$ 时，$\Delta_\text{r} H_\text{m} = -50.0\ \text{kJ}\cdot\text{mol}^{-1}$，电池反应平衡常数 $K^\ominus = 0.0301$，$E^\ominus(\ce{Ag+ / Ag}) = 0.800\ \text{V}$，设活度系数均为 1。
\begin{enumerate}[label=（\arabic*）, itemsep=1ex]
    \item 写出电极与电池反应；
    \item 计算 298 K 时 $\ce{AgBr(s)}$ 的 $K_\text{sp}$；
    \item 求电池反应的可逆反应热 $Q_R$；
    \item 计算电池的温度系数。
\end{enumerate}

\vspace{0.5cm}
\noindent\textbf{八、计算题(13 pts.)}

\ce{NO} 高温均相分解是二级反应，反应为：$\ce{2NO(g) -> N2(g) + O2(g)}$。实验测得 1423 K 时速率常数为 $1.843 \times 10^{-3}\ \text{dm}^3\cdot\text{mol}^{-1}\cdot\text{s}^{-1}$，1681 K 时速率常数为 $5.743 \times 10^{-3}\ \text{dm}^3\cdot\text{mol}^{-1}\cdot\text{s}^{-1}$。求：
\begin{enumerate}[label=（\arabic*）, itemsep=1ex]
    \item 等容反应的活化能；
    \item 用过渡态理论计算 1423 K 时反应的 $\Delta^\ddagger H_\text{m}^\ominus$、$\Delta^\ddagger U_\text{m}^\ominus$、$\Delta^\ddagger S_\text{m}^\ominus\ (c^\ominus)$ 及 $\Delta^\ddagger G_\text{m}^\ominus\ (c^\ominus)$。
\end{enumerate}